To rigorously prove that the extremal equality $|\{p : \sigma(p) \ne 0\}| = n$ cannot hold for any strictly convex polygon $P$ with at least 3 distinct edge-direction classes, we track the alternating sign structures and identify strictly uncancelled points inside the Minkowski sum $Q+P$, where $Q = \operatorname{conv}(T)$.

### 1. Geometric Setup and Cycle Constraint
The set $T$ consists of translates $t_i$ exposed at each normal cone of $s_i$. $T$ traces the boundary of the convex polygon $Q = \operatorname{conv}(T)$. 
For each edge $[s_i, s_{i+1}]$ of $P$ with direction $a_i = s_{i+1} - s_i$, the translates $F_i = T \cap [t_i, t_{i+1}]$ form an arithmetic progression of step $a_i$ and length $m_i$. Thus, $t_{i+1} = t_i + m_i a_i$. 
Traversing the full boundary of $Q$ implies:
$$ \sum_{i=1}^n m_i a_i = 0 $$
The points in $F_i$ have alternating signs. For consistency, $\sigma(t_{i+1}) = (-1)^{m_i} \sigma(t_i)$. Tracing the full cycle dictates $\prod_{i=1}^n (-1)^{m_i} = 1$, making $\sum_{i=1}^n m_i$ an even integer. 

We write the convolution formally in the group ring $\mathbb{Z}[\mathbb{R}^2]$ as $T \cdot S = \sum_{p} \sigma(p) [p]$.
Let $C_i = \epsilon_i \sum_{k=0}^{m_i-1} (-1)^k[t_i + k a_i]$ where $\epsilon_i = \sigma(t_i)$. It is straightforward to show that the points of $T \cdot S$ on the boundary of $Q+P$ perfectly telescope, leaving exactly the $n$ "lonely vertices" $V_i = t_i + s_i$, each with coefficient $\epsilon_i \in \{+1, -1\}$.
For the extremal equality to hold, **every other point strictly inside $Q+P$ must cancel to $0$**.

### 2. Identifying Extra Uncancelled Points
We classify the shape of $Q$ based on its positive coefficients $m_i > 0$. Because $\sum m_i a_i = 0$ and $P$ has at least 3 distinct direction classes, we encounter three mutually exhaustive geometric cases:

**Case A: $Q$ has two adjacent positive edges.**
Suppose $m_i > 0$ and $m_{i+1} > 0$. Consider the target point:
$$ P^* = t_{i+1} - a_i + s_{i+2} = t_{i+1} + a_{i+1} + s_i $$
This point lies strictly inside $Q+P$. Let us find all pairs $(t,s) \in T \times S$ that sum to $P^*$.
Because $P^*$ is near the vertex $V_{i+1} = t_{i+1} + s_{i+1}$, the extreme convex geometry restricts the components:
1. $t = t_{i+1} - a_i \in F_i$ and $s = s_{i+2}$. Its sign in $T$ is $-\epsilon_{i+1}$.
2. $t = t_{i+1} + a_{i+1} \in F_{i+1}$ and $s = s_i$. Its sign in $T$ is $-\epsilon_{i+1}$.

No other combinations in $T \times S$ can reach $P^*$. Thus, the coefficient of $P^*$ in the convolution is $(-\epsilon_{i+1}) + (-\epsilon_{i+1}) = -2\epsilon_{i+1} \neq 0$. This gives an extra uncancelled point not among the $n$ lonely vertices, breaking the extremal equality.

**Case B: $Q$ has non-adjacent positive edges separated by $m_j=0$, but without opposite parallel directions.**
To avoid Case A, positive $m_i$ edges must be separated by edges with $m_j = 0$. Suppose $m_i > 0$ and $m_k > 0$, and $m_j = 0$ for all $i < j < k$. 
Because $m_j = 0$, the corresponding translates are stationary: $t_{i+1} = t_{i+2} = \dots = t_k$. Let this shared vertex of $Q$ be $v$.
Consider the target point:
$$ P^{**} = v - a_i + s_k $$
This point is formed by $(v - a_i) \in F_i$ paired with $s_k$, generating the sign $-\epsilon_{i+1}$.
For $P^{**}$ to cancel, it must be matched by another pair. The only viable vertex in $S$ near this location is $s_{k+1}$, demanding $t \in T$ such that $t + s_{k+1} = v - a_i + s_k$. This implies $t = v - a_i - a_k$. Since $v$ is an extreme point in $T$, this is only possible if $t = v$ (which belongs to $F_k$ with sign $\epsilon_k = \epsilon_{i+1}$), severely forcing $a_k = -a_i$.
If $a_k \neq -a_i$, the point $P^{**}$ is uniquely formed and strictly uncancelled, breaking the equality.

**Case C: $Q$ uses only opposite parallel directions (Line Segment Degeneracy).**
To avoid both $P^*$ and $P^{**}$, $Q$ must only have positive lengths for a single pair of strictly opposite parallel edges. Hence $a_k = -a_i$ (so $a_i$ and $a_k$ share a direction class, satisfying $a_k = -\lambda a_i$).
Consequently, $m_j = 0$ for all other indices, meaning $Q$ degenerates into a 1-dimensional line segment. 
The chains $F_i$ and $F_k$ act precisely on this single segment. As warned by the hint, since $\sum m_x a_x = 0$, they must span the exact same geometric length back and forth. Because $T$ is a proper set (no repeating elements), the only valid layout is a single alternating 1D chain of step $a_i$.

However, since $P$ possesses **at least 3 distinct edge-direction classes**, $S$ has vertices (say $s_r$) that are not on the edges parallel to $a_i$.
When we convolve the 1D segment $T$ with such a vertex $s_r$, we produce an entire line segment of points $T + s_r$. 
For any of these points to be cancelled, there would have to exist another vertex $s_l \in S$ such that the segment $T + s_l$ perfectly overlaps with $T + s_r$, requiring the chord $s_l - s_r$ to be parallel to $a_i$. Since $P$ is strictly convex, chords parallel to an edge can only be the opposing edge itself. But we've established $P$ possesses classes beyond this pair!
Therefore, the points generated by $T + s_r$ face no possible cancellation, flooding the inside of $Q+P$ with at least $m_i + 1 \ge 2$ extra non-zero points per non-parallel vertex.

### Conclusion
In every case, the presence of at least three edge-direction classes forces EITHER a structural overlap that violates $T$ being a set, OR strictly creates extra uncancelled points (such as $P^*$ or $P^{**}$) trapped inside $Q+P$. Consequently, $|\{p : \sigma(p) \neq 0\}|$ strictly exceeds $n$, making the extremal equality impossible for all triangles, non-parallelogram quadrilaterals, and all $n$-gons where $n \ge 5$.